\chapter{No Singularities in Nature}
\label{chap:bounce}

A singularity is not a feature of the world; it is the noise a theory makes when
it has been pushed past what it knows. General relativity, left to itself,
predicts that gravitational collapse ends in a point of infinite density---and
then it stops being able to say anything at all. The founding wager of this book
(Chapter~\ref{chap:axioms}) is to take that breakdown at face value: physics is
continuous and conservative, so the infinity is an artifact, and \emph{something}
must intervene before it is reached. In Einstein--Cartan gravity, something does.
The collapse does not terminate. It \emph{bounces}. Everything else in this book
is the working-out of that single refusal to admit a singularity.

\section{What stops a collapse}

Einstein--Cartan gravity is general relativity with one restriction lifted: the
connection is allowed an antisymmetric part, torsion, which couples to the spin
of fermions (Axiom~\ref{ax:ec}). That spin-torsion coupling turning gravitational
collapse into a non-singular bounce is the mechanism developed in a series of
papers by \citet{poplawski2010,poplawski2012}. At everyday densities this changes
nothing measurable. But it adds a second term to the high-density stress--energy, and the
two terms scale differently. The \emph{gravitational} term is attractive and
grows linearly with mass--energy density,
\[
  T_{\mathrm{grav}} \sim \rho c^2 ,
\]
while the \emph{torsion} term is repulsive and grows as the \emph{square} of the
spin density $s$,
\[
  T_{\mathrm{tors}} \sim \frac{G\hbar^2 s^2}{c^4} .
\]
At nuclear density the ratio $T_{\mathrm{tors}}/T_{\mathrm{grav}}$ is about
$10^{-40}$---utterly negligible, which is why no laboratory has ever seen it. But
compression packs more fermions into each volume, each carrying spin $\hbar/2$,
and the quadratic term climbs faster than the linear one. Squeeze hard enough and
repulsion overtakes attraction. The inward fall is halted from within, by the
spin of the matter itself, and reverses. No outside agent, no new field, no
fine-tuning: just a term that general relativity left out because torsion was
assumed to vanish.

\section{The bounce density}
\label{sec:bounce-density}

Setting the two terms equal fixes the density at which the reversal happens. For
ordinary fermionic matter, with about one unit of spin $\hbar$ per nucleon mass,
\[
  \rhoc \;\sim\; \frac{m_n^2 c^4}{G\hbar^4}
  \;\sim\; 10^{50}\,\mathrm{kg\,m^{-3}} .
\]
This sits 33 orders of magnitude above nuclear density and 46 below the Planck
density---comfortably in a regime where Einstein--Cartan is still classical and
trustworthy, far from where quantum gravity proper would be needed. The exact
figure depends on the fermion content and on coefficients in the
Einstein--Cartan Lagrangian, and pinning it precisely is one of the calculation
targets of the simulation program (Chapter~\ref{chap:sims}).

\begin{figure}[t]
  \centering
  \includegraphics[width=\textwidth]{bounce}
  \caption{The Einstein--Cartan torsion bounce (a 1D spherical collapse, radiation
    fluid, in dimensionless units $8\pi G/3=\rhoc=1$). \emph{Left:} the collapsing
    universe reaches a minimum radius $a_{\min}$ and re-expands (blue) exactly where
    general relativity would run to a singularity, $a\to0$ (red dashed).
    \emph{Right:} the density saturates at the bounce density $\rhoc$ and the Hubble
    rate crosses zero---collapse turning into expansion, with no singularity ever
    reached. Reproducible as \texttt{figures/scripts/ch09\_bounce.py}; the full
    numerical treatment is Chapter~\ref{chap:sims}.}
  \label{fig:bounce}
\end{figure}

\paragraph{The shape of the turnaround.} It helps to see the bounce as a particle
rolling in a potential. Writing the effective Friedmann equation
$H^2=a^{-n}(1-a^{-n})$ (dimensionless, with $\rhoc$ and $a_{\min}$ set to one and
$n=3(1+w)$) as a zero-energy mechanics problem $\dot a^2+V(a)=0$ gives
\begin{equation}
  V(a)=a^{2-2n}-a^{2-n},
  \label{eq:bounce-potential}
\end{equation}
a \emph{single repulsive wall} (\Cref{fig:bouncepot}). The positive term
$a^{2-2n}$ is the torsion contribution---the $-\rho^2/\rhoc$ spin term---and it
walls off small $a$; the negative term $-a^{2-n}$ is the ordinary matter pull. They
cross at $a_{\min}$, where the contracting universe turns around ($V=0$, slope
$-n$). There is no second minimum and nothing tuned: the only things that fix the
shape are the equation of state $w$ (through $n$) and the density scale $\rhoc$. The
microphysics that fixes \emph{that} scale---the strong four-fermion coupling at the
wall---is where the framework's deepest open question lives (Appendix~\ref{app:openq}).

\begin{figure}[t]
  \centering
  \includegraphics[width=0.82\textwidth]{bounce_potential}
  \caption{The bounce as a single repulsive wall, \eqref{eq:bounce-potential}. The
    torsion term $a^{2-2n}$ (dotted) forbids small $a$; the matter pull $-a^{2-n}$
    (dashed) dominates at large $a$; they meet at $a_{\min}$ ($V=0$), the turnaround.
    The classically forbidden region is $V>0$ ($a<a_{\min}$); the expanding universe
    lives in $V<0$. The only knobs are $w$ and $\rhoc$. Reproducible as
    \texttt{figures/scripts/ch02\_bounce\_potential.py}.}
  \label{fig:bouncepot}
\end{figure}

\paragraph{A name.} The density at which Cartan's torsion forces the turnaround is
where everything is dissolved and then extruded clean---a \emph{katharsis}. Cartan
$+$ \emph{katharsis} gives \emph{Cartasis}: one may call $\rhoc$ the Cartasis density
and the bounce surface the Cartasis membrane. Elsewhere they are simply the
\emph{bounce density} and the \emph{bounce membrane}.

\paragraph{What is actually there.} At $\rhoc$ the familiar particle picture is
gone. The matter is compressed past quark deconfinement, past electroweak
restoration, past any regime in which the Standard Model's quasiparticles are
even defined. What remains is continuous field content---energy--momentum
density, spin density, chiral current density---and what crosses the membrane is
therefore not \emph{particles} but \emph{conserved currents}: total
energy--momentum, total spin, total chiral charge, and (subject to electroweak
anomaly effects) baryon and lepton number. Nothing structured survives the
crossing. Whatever a world is built from on the far side must be rebuilt from
elementary field content---a point that will matter both for what a universe
inherits (Chapter~\ref{chap:dawn}) and for why a stray ordered object cannot
simply pass through (Chapter~\ref{chap:void}).

\paragraph{Why this is gravity's coupling, and a neutron star proves it.} The repulsion
that halts the collapse is a spin--spin (four-fermion) interaction, and the density at
which it wins scales \emph{inversely} with that interaction's strength: a stronger
coupling bounces sooner, at lower density. The coupling here is gravitational---$\kappa
\sim G/c^4$, Planck-suppressed---which is precisely why $\rhoc$ is so enormous. It is worth
checking that it could not instead be the much stronger coupling that, in Part~II, gives
particles their mass (a strong condensate at $\Lambda\sim 246~\mathrm{GeV}$, some $(M_{\mathrm{Pl}}/\Lambda)^2\sim10^{33}$ times stronger than gravity). If \emph{that} coupling drove the
bounce, the inverse scaling would drop $\rhoc$ by the same $10^{33}$---to about
$4\times10^{16}~\mathrm{kg\,m^{-3}}$, a fifth of nuclear density. But neutron-star cores
reach several times nuclear density and manifestly do \emph{not} bounce. So the strong,
mass-giving coupling is excluded as the bounce driver by the mere existence of neutron
stars: the bounce is gravitational, and the mass sector is a separate, far stronger force.
This is not a worry to defer but a check one can run on a napkin, and it has a clean
consequence: because the bounce temperature $T\sim\rhoc^{1/4}\sim10^{7}~\mathrm{GeV}$ is some
$10^{5}$ times the condensation scale $\Lambda$, the mass condensate is fully melted at the
bounce (consistent with the restored, massless matter above), and the mass sector contributes
there only as extra radiation---a revision to the degree-of-freedom count, not to the
gravitational repulsion. The bounce of Part~I therefore survives the propagating-torsion
extension that Part~II's mass scale requires (Chapter~\ref{chap:weltformel};
Figure~\ref{fig:bounceconsistency}).

\begin{figure}[t]
  \centering
  \includegraphics[width=\textwidth]{bounce_consistency}
  \caption{The bounce coupling is gravitational, not the mass coupling---and a neutron star
  decides it. The bounce density scales inversely with the four-fermion coupling: the
  gravitational $G/c^4$ puts it at $\rhoc\sim10^{50}~\mathrm{kg\,m^{-3}}$, whereas the
  $\sim10^{33}$-times-stronger mass-giving coupling would put it at $\sim0.2\times$ nuclear
  density---below the neutron-star cores that exist without bouncing. The two scales are
  necessarily two couplings. Reproducible from
  \texttt{sims/src/cartasis\_sims/bounce\_consistency.py}.}
  \label{fig:bounceconsistency}
\end{figure}

\section{The geometry of the bounce}

\subsection{Radius becomes time}

Inside a Schwarzschild horizon the radial coordinate $r$ is timelike: future-%
directed motion is motion inward, and this is ordinary general relativity. The
bounce uses that fact rather than fighting it. Torsion pushes the matter
``outward'' in $r$---but since $r$ is timelike there, outward in $r$ is not a
path back through the horizon; it is a path into the \emph{future} of the bounce.
That future, read in the new region's own coordinates, is the cosmological future
of a fresh spacetime. The parent's radial direction $r$ becomes, smoothly across
the membrane, the new world's time direction $\tau$.

\subsection{Continuity across the membrane}

Because nothing is allowed to be discontinuous, the metric and connection must
match across the membrane. The induced metric agrees from both sides; the
extrinsic curvature matches under junction conditions of Israel--Darmois type,
modified to carry torsion; and stress--energy and spin current flow through
unbroken. These conditions are simply the precise statement of ``the bounce is
continuous''---a system of equations relating the parent-side and child-side
fields at the surface, and the thing a faithful simulation must reproduce.

\subsection{A persistent interface, not a single instant}

A black hole that keeps accreting keeps feeding the bounce. New matter falls in,
compresses, reaches the bounce density, and joins the membrane---so the surface
is not a one-time event but a persistent interface that lasts as long as the
parent hole does. This is the framework's deepest departure from the standard
picture, and it deserves to be stated plainly: a child universe's Big Bang is
\emph{not} an instant in its past but a surface in its present, still dawning.
Matter is injected across it throughout the parent's life, not only at the first
moment---the beginning never ends, it merely recedes to the horizon. That ongoing
injection is precisely what the child's inhabitants will later measure as their
dark sector (Chapter~\ref{chap:home}), so the dark sky is the standing evidence
that the dawn is still in progress. A dawn closes only when its \emph{feed} runs dry. For a
child that is the parent ceasing to accrete, or evaporating; but an original-generation
universe has no parent and instead gorges directly on the void, so \emph{its} dawn is ongoing
too---for as long as it draws matter from the vacuum. Only an old, fed-out OGU has a beginning
that is truly finished, and it then evaporates in its turn, after a span beyond
$10^{100}$ years. The OGU is therefore not the exception to the ongoing dawn but a different
\emph{source} for it---vacuum rather than parent. What it lacks is not the dawn but the dark
sector: a diffuse void-feed presses no concentrated parent mass behind a membrane, so an OGU
dawns without the dark matter and dark energy that betray a child's ongoing injection
(Chapter~\ref{chap:home}).

\subsection{The parent's horizon is the child's first light}

The same surface wears three faces. From the parent's exterior it sits just
inside the event horizon. From the parent's interior it is wherever infalling
matter is currently reaching the bounce density. From the child's interior,
looking as far back in time as it can, it is the boundary of the past at maximum
lookback---the child's Big-Bang surface. The parent's horizon region and the
child's earliest sky are one surface seen from opposite sides. This single
identification is what later lets the cosmic microwave background be read as the
parent's Hawking glow (Chapter~\ref{chap:home}).

\section{What the bounce does \emph{not} do}
\label{sec:cpt-even}

The bounce reliably turns a black hole into an expanding world, but it does not
also \emph{make} the matter that fills it---it is no mirror that turns infalling
antimatter into matter and hands us the cosmic excess for free. The interaction
that drives it forbids that, and the proof is short. When torsion is integrated
out of Einstein--Cartan--Sciama--Kibble gravity, the Dirac sector is left with a
single four-fermion contact term, the Hehl--Datta interaction~\citep{hehl1976},
\[
  \Delta\mathcal{L}_{\mathrm{tors}}
  = -\,\frac{3}{16}\,\frac{8\pi G}{c^4}\,
    \big(\bar\psi\gamma^5\gamma_a\psi\big)\big(\bar\psi\gamma^5\gamma^a\psi\big),
\]
the negative square of the axial current. This is the very term whose energy
density grows as the square of the spin density and overwhelms gravity at $\rhoc$:
it \emph{is} the bounce, so the bounce inherits its discrete symmetries. The axial
current is even under charge conjugation, its contraction is even under parity,
and the squared bilinear is even under time reversal. The interaction that drives
the bounce is thus separately $C$-, $P$-, and $T$-even, and three things follow
that no amount of modelling can wish away:

\begin{enumerate}
  \item \textbf{No flip.} A $C$-even interaction cannot turn matter into
    antimatter. The membrane does not charge-conjugate what passes through it.
  \item \textbf{No filter from nothing.} A $C$- and $P$-even interaction cannot
    \emph{create} a chiral imbalance: zero net axial density in gives zero out.
  \item \textbf{No arrow reversal.} A $T$-even dynamics does not run proper time
    backwards.
\end{enumerate}

The bounce is real, robust, and generous---it reliably turns a black hole into an
expanding world---but it is not a baryogenesis engine. Any matter--antimatter
asymmetry must be \emph{supplied}, not manufactured at the membrane. Where it
actually comes from is the business of Chapter~\ref{chap:dawn}.

\section{The arrow is struck at the bounce: a Janus point}
\label{sec:janus}

If the bounce does not reverse time, where does a universe's arrow of time come
from? It is \emph{made} at the membrane. Along the matter's worldlines proper
time runs monotonically; nothing turns around, the scale factor merely passes
through a minimum. What turns around is the \emph{thermodynamic} arrow. Infalling
parent material arrives from a structured, high-entropy universe; compression
toward $\rhoc$ dismantles that structure and smooths the geometry---the Weyl
curvature is driven toward zero---so gravitational entropy \emph{falls} toward the
membrane. On the far side, expansion lets structure, and eventually black holes,
form again, so entropy \emph{rises} away from it. Entropy is therefore at a
\emph{minimum at the membrane}, and the locally defined arrow points away from the
bounce on both sides at once: a \emph{Janus point}~\citep{barbour2014}, one
spacetime with a single monotonic proper time and two thermodynamic arrows
diverging from a central low.

A newborn universe thus inherits a genuine low-entropy past for free, with no time
flip required---the long-postulated ``Past Hypothesis'' arriving as a consequence
of bounce geometry rather than as an unexplained initial condition. (That this
low-entropy past is \emph{forced}, and what it means for the timeless void from
which the very first bounce is struck, is taken up in Chapter~\ref{chap:void}; how
the arrow ends up locked to the matter--antimatter choice is taken up in
Chapter~\ref{chap:dawn}.) The old slogan that ``$r$ becomes timelike, so time
inverts'' is a statement about Schwarzschild coordinates inside a horizon, not a
physical reversal of anyone's clock.

\section{What forms after the bounce}

On the far side of the membrane the bounce dynamics open a new spacetime
region---a child universe with its own time direction and its own spatial
extent, the baby-universe scenario of \citet{poplawski2016} and, earlier, of
\citet{frolov1990}---and from inside
it looks like a perfectly ordinary hot Big Bang: it
begins at maximum density, expands fast, cools, passes through radiation- and
matter-dominated eras, and grows large-scale structure. The early, very rapid
expansion is driven by the torsion repulsion at its post-bounce peak and fades as
the density drops, giving an inflation-like epoch with no separate
inflaton~\citep{linde1982}; its
quantitative shape, and the perturbation spectrum it should leave, are developed
in Chapter~\ref{chap:dawn} and put to the test in Chapter~\ref{chap:nowobs}.

From the parent's side none of this is visible as a spatial object; the parent
simply sees a black hole where the collapse occurred. The child's vast interior
lies in the future of the bounce, in directions orthogonal to the parent's
geometry---which produces the framework's most cheerfully counterintuitive
feature, a \emph{Tardis}: a universe far larger inside than its parent's horizon
would suggest. Our own matter content, some $10^{53}\,\mathrm{kg}$, sits at the
bounce density in about $10^{3}\,\mathrm{m^3}$---less than half an Olympic
pool---yet from inside we measure a universe tens of billions of light-years
across. There is no contradiction: the small volume is the matter at the bounce
moment in the parent's coordinates, the large one is the cosmological extent after
expansion in the child's, and they are simply different quantities. Even plain
general relativity already has the spatial volume inside a horizon growing without
bound~\citep{christodoulou2015}; the bounce reinterprets that growth as the
expansion of a child universe rather than the approach to a singularity we have
agreed nature does not permit.
